\section{Hold'em Games}
In diesem Kapitel stelle ich vor, welche Pokervarianten im Rahmen dieser Arbeit mittels CFR gelöst werden sollen. 
Hierbei erkläre ich zunächst grundlegend wodurch sich die Untergruppierung von Pokerspielen, Hold'em-Spiele, auszeichnen und stelle dann einige bekannte Varianten vor.
Denn, wenn, man die Grundstruktur von Hold'em-Spielen definiert hat, ist es einfach diese Struktur zu erweitern und alle folgenden Poker Hold'em Varianten darzustellen.
Anschließend beschreibe ich die Varianten welche in dieser Arbeit gelöst werden und warum ich diese Auswahl getroffen habe. 

Ich beschäftige mich in dieser Arbeit ausschließlich mit Heads-Up Limit Versionen von Hold'em-Spielen.
Heads-Up bedeutet, dass das Spiel nur zwei Teilnehmer beinhaltet. 
Die Klassifikation Limit bezieht sich darauf, dass die Setzgrößen der Spieler auf vorher festgelegte Beträge beschränkt sind.
Diese zwei Einschränkungen reduzieren die Komplexität enorm. 
Die folgende Definition gilt für jegliche Heads-Up Limit Hold'em-Spiele.

Eine Iteration eines Hold'em-Spiels nennt man eine Hand.
Eine Hand widerum besteht aus mehreren Runden.
In der ersten Runde wird eine bestimmte Anzahl an privaten Karten (Hole Cards) an die zwei Teilnehmer verteilt.
In \emph{jeder} Runde wird eine bestimmte Anzahl (0 auch möglich), an öffentlichen Karten (Board Cards) aufgedeckt. 

Nachdem Karten ausgeteilt oder aufgedeckt wurden, folgt immer eine Setz Runde, in welcher die Spieler abwechselnd Entscheidungen treffen.
Ihnen stehen grundlegend 3 Auswahlmöglichkeiten zur Verfügung: \grqq fold, call und raise\glqq.
Die Begrifflichkeiten können sich abhängig von der Quelle und dem Spiel etwas unterscheiden.
\\
Wählt ein Spieler die Aktion \emph{fold}, endet die Hand sofort und der andere Spieler gewinnt den Pot.
Der Pot ist die Ansammlung aller Einsätze die im Verlauf der Hand gesetzt wurden.
\\
Bei einem \emph{call}, setzt er so viele Chips in den Pot, bis er genauso viele Chips in den Pot eingezahlt hat wie der andere Spieler.
Die Anzahl an Chips kann auch 0 betragen. 
In Pokervarianten welche zu Unterhaltungs und nicht zu Forschungszwecken gespielt werden, nennt man das dann \emph{check}. 
\\
Bei der Aktion \emph{raise}, setzt der Spieler zunächst genauso viel Chips in den Pot wie sein Gegner, analog zum \emph{call}.
Anschließend setzt er dann zusätzlich noch einen durch die Spielregeln festgelegten Betrag an Chips. 

Die Spieler treffen solange abwechselnd Entscheidungen bis entweder ein Spieler folded, in diesem Fall endet die Hand sofort, oder bis ein Spieler called (solange es nicht die erste Aktion der Runde ist).
In dem Fall das ein Spieler called, endet nur die Runde und nicht die Hand.
Um eine endlose Abfolge aus Raises zu vermeiden, gibt es eine Obergrenze für die Anzahl an Raises nacheinander ausgeführt werden dürfen.
Wenn das Setzlimit erreicht ist, bleiben dem Spieler welcher am Zug ist, nur die Optionen Call und Fold.
Wenn bis zum Ende der finalen Setzrunde kein Spieler gefolded hat, kommt es zu einem \emph{Showdown}.
Bei einem Showdown offenbaren beide Spieler ihre Hole Cards und bilden aus allen zu Verfügung stehenden Karten (Hole + Board Cards) die bestmögliche Pokerhand.
Der Spieler mit der stärkeren Hand gewinnt den Pot, welche möglichen Pokerhände es gibt, und welcher Rangordnung diese folgen, hängt von der Variante des Spiels ab.
Eine Übersicht über die verschiedenen Pokerhände und deren Bedeutung findet sich unter \url{https://en.wikipedia.org/wiki/List_of_poker_hands}.
\cite{southey2005bayes}

Da wir jetzt eine grundlegende Definition eines Heads-Up Limit Holdem Spiels haben, können wir nun, durch anpassen einzelner Parameter jegliche Varianten definieren.
Die relevantesten Aspekte sind: Die Anzahl der Hole und Board Cards, die Anzahl der Setzrunden, dem Setzlimit pro Runde, den im Deck enthaltenden Karten, die möglichen Pokerhände, die festgelegten Setzgrößen und noch einiges mehr.


\subsection{Kuhn Poker}
Die simpelste definierte Pokervariante ist One Card Poker auch Kuhn Poker gennant.
Der Name Kuhn Poker stammt durch den Erfinder dieses Spiels Harald Kuhn, welcher es in seinem Paper aus dem Jahr 1951 mit dem Titel \grqq Simplified Two-Person Poker \glqq beschrieb. \\
Kuhn Poker wird mit einem Kartendeck bestehend aus 3 Karten gespielt.
In seinem Paper sind die Karten mit 1,2 und 3 bennant, wobei 1 die schlechteste und 3 die beste Karte ist.
Für diese Arbeit wird die Notation J, Q, K verwendet, die den Rängen eines Standard-Kartenspiels entspricht.

Zu Beginn zahlen beide Spieler einen Pflichteinsatz in Höhe eines Chips in den Pot, dies wird auch Ante genannt. 
Beide Spieler erhalten eine private Karte und es werden keine öffentlichen Karten ausgeteilt.
Nach dem austeilen der privaten Karte startet die erste und einzige Setzrunde, in welcher die Spieler jeweils einen Chip in den Pot setzen dürfen, womit das Setzlimit bei Eins liegt.
Im Falle eines Showdowns, gewinnt die höhere Karte $J < Q < K$.

Kuhn definiert in seinem Paper auch noch andere Versionen seines Spiels indem er das Verhältnis zwischen Ante und Setzgröße anpasst.
Der eben beschriebene Fall in welchem der Ante genauso groß ist wie die Setzgröße, wird von ihm Fall 2 genannt.
Fall 1 liegt vor, wenn die Setzgröße kleiner als der Ante ist.
Fall 3 tritt ein, wenn der Ante kleiner als die Setzgröße ist und die Setzgröße kleiner als das Doppelte des Antes ist.
Fall 4 liegt vor, wenn die Setzgröße dem Doppelten des Antes entspricht.
\cite{kuhn1951simplified}
\subsection{Leduc Hold'em}
Leduc Hold'em~\cite{southey2005bayes} ist eine vereinfachte Hold'em Variante für zwei Spieler.

Leduc Hold'em wird mit einem Kartendeck bestehend aus zwei Suits mit drei Karten pro Suit gespielt, insgesamt sechs Karten.
Die Karten Notation verwendet Rank und Suit, beispielsweise Js, Jh, Qs, Qh, Ks, Kh.

Das Spiel besteht aus zwei Runden.
In der ersten Runde erhält jeder Spieler eine private Karte.
In der zweiten Runde wird eine öffentliche Board Karte aufgedeckt.

Zu Beginn zahlen beide Spieler einen Ante von einem Chip in den Pot.
Es gibt ein Maximum von zwei Bets pro Runde.
Die Setzgrößen betragen 2 Chips in der ersten Runde und 4 Chips in der zweiten Runde.

Die möglichen Pokerhände sind entweder ein Paar oder nur eine High Card. 
Ein paar gewinnt gegen eine High Card.


\subsection{Zwölf Karten Poker}
Die folgende Pokervariante ist von mir selbst entworfen worden, als eine Zwischenstufe der Komplexität zwischen Leduc Hold'em und Rhode Island Hold'em.
Die 12 Karten Variante verwendet ein Deck mit zwölf Karten: J, Q, K und A in drei Suits (s, h, d).

Es gibt drei Runden.
In der ersten Runde erhält jeder Spieler eine private Karte.
In der zweiten und dritten Runde wird jeweils eine öffentliche Board Karte aufgedeckt.

Es gibt ein Maximum von zwei Bets pro Runde.
Die Setzgrößen betragen 2 Chips in der ersten Runde, 4 Chips in der zweiten Runde und 8 Chips in der dritten Runde.
Zu Beginn zahlen beide Spieler einen Ante von 1 Chip in den Pot.

Das Spiel endet nach der dritten Runde mit einem Showdown.
Die Hand besteht aus der Kombination der privaten Karte des Spielers und den beiden öffentlichen Board Karten.
Die Hand Rankings sind: Three of a Kind schlägt ein Paar, ein Paar schlägt High Card.



\subsection{Rhode Island Hold'em}
Rhode Island Hold'em~\cite{shi2001rhode} wurde von Shi und Littman im Jahr 2001 als Testbed für AI Forschung entwickelt.

Rhode Island Hold'em verwendet ein Standard Kartendeck mit 52 Karten.
Zu Beginn zahlen beide Spieler einen Ante von 5 Chips in den Pot.
Jeder Spieler erhält eine private Karte.

Das Spiel besteht aus drei Runden.
Die erste Runde findet nach dem Austeilen der privaten Karten statt.
Die Bet Size beträgt 10 Chips und es gibt ein Limit von drei Raises pro Runde.
Nach der ersten Runde wird eine öffentliche Board Karte, der Flop, aufgedeckt.
Die zweite Runde findet nach dem Flop statt.
Die Bet Size beträgt 20 Chips.
Nach der zweiten Runde wird eine weitere öffentliche Board Karte, der Turn, aufgedeckt.
Die dritte und finale Runde findet nach dem Turn statt.
Die Bet Size beträgt ebenfalls 20 Chips.

Beim Showdown bildet jeder Spieler die beste 3 Card Poker Hand aus seiner privaten Karte und den beiden öffentlichen Board Karten.
Die Hand Rankings für 3 Card Poker unterscheiden sich von 5 Card Poker.
Three of a Kind schlägt Straight und Flush.
Straight schlägt Flush.
Flush schlägt Pair und Pair schlägt High Card.
Bei einem Unentschieden wird der Pot gleichmäßig geteilt.


\subsection{Limit Texas Hold'em}
Texas Hold'em~\cite{southey2005bayes} ist das häufigste Format für Hold'em.

Texas Hold'em verwendet ein Standard Kartendeck mit 52 Karten, bestehend aus 2 bis Ass in vier Suits.

Das Spiel besteht aus vier Betting Runden.
In der ersten Runde erhalten die Spieler jeweils zwei private Karten.
In der zweiten Runde, dem Flop, werden drei Board Karten aufgedeckt.
In der dritten Runde, dem Turn, wird eine vierte Board Karte aufgedeckt.
In der vierten Runde, dem River, wird eine fünfte Board Karte aufgedeckt.

Es gibt ein Maximum von vier Bets pro Runde.
Die festgelegten Bet Sizes betragen 10 Chips in den ersten beiden Runden und 20 Chips in den letzten beiden Runden.
In diesem Variation beginnt das Spiel nicht mit einem Ante, bei dem beide Spieler die gleiche Anzahl an Chips in den Pot zahlen, sondern mit einer Struktur die man Blinds nennt.
Dabei zahlt der erste Spieler 5 Chips in den Pot und der zweite 10 Chips.

Die möglichen Pokerhände sind Royal Flush, Straight Flush, Four of a Kind, Full House, Flush, Straight, Three of a Kind, Two Pair, One Pair und High Card.
Diese Vielfalt an möglichen Händen trägt zur hohen Komplexität des Spiels bei.



